\documentclass[10pt,a4paper]{article}
\usepackage[T1]{fontenc}
\usepackage[utf8]{inputenc}
\usepackage[ngerman]{babel}
\usepackage[a4paper,margin=12mm]{geometry}
\usepackage{lmodern}
\usepackage[table]{xcolor}
\usepackage{tabularx,array,booktabs,amsmath,microtype}
\usepackage[most]{tcolorbox}
\usepackage{tikz}
\usetikzlibrary{arrows.meta,calc}
\pagestyle{empty}
\setlength{\parindent}{0pt}
\setlength{\parskip}{0pt}
\renewcommand{\familydefault}{\sfdefault}
\renewcommand{\arraystretch}{1.12}
\definecolor{Navy}{HTML}{17324D}
\definecolor{Blue}{HTML}{3B6EA5}
\definecolor{Sky}{HTML}{EAF3FA}
\definecolor{Mint}{HTML}{E8F5EF}
\definecolor{Gold}{HTML}{F3B33D}
\definecolor{Ink}{HTML}{1E2935}
\definecolor{Muted}{HTML}{617080}
\definecolor{Line}{HTML}{D8E1E8}
\color{Ink}
\newcommand{\unit}[1]{\,\mathrm{#1}}
\newcommand{\unitfrac}[2]{\,\frac{\mathrm{#1}}{\mathrm{#2}}}
\newcommand{\sectiontitle}[1]{\par\vspace{2mm}{\color{Blue}\bfseries\large #1}\par\vspace{.3mm}{\color{Blue}\rule{\linewidth}{.7pt}}\par\vspace{1mm}}
\newcommand{\step}[2]{\par\vspace{1.8mm}\textbf{\textcolor{Blue}{#1}\quad #2}\par\vspace{.5mm}}
\newcommand{\sheethead}[1]{\par\color{Ink}\leavevmode%
  \colorbox{Navy}{\parbox{\dimexpr\linewidth-2\fboxsep\relax}{\vspace{1.3mm}\color{white}\small\bfseries PHYSIK \textperiodcentered{} KLASSE 11\hfill\textcolor{Gold}{LÖSUNGEN}\par\vspace{1.1mm}{\large Lösungsblatt \textperiodcentered{} Schritt für Schritt}\vspace{1.1mm}}}
  \par\vspace{1.6mm}{\color{Navy}\fontsize{15}{18}\selectfont\bfseries Übungen zur Zentripetalkraft\par}
  \vspace{.5mm}\textbf{Aufgabe 7}\hfill\textcolor{Muted}{#1}\par\vspace{1mm}}
\newcommand{\pagefooter}[1]{\vfill{\color{Line}\rule{\linewidth}{.5pt}}\par\vspace{1mm}\begin{tabularx}{\linewidth}{@{}Xc r@{}}{\leavevmode\small\color{Muted}Klasse 11 \textperiodcentered{} Physik \textperiodcentered{} Kreisbewegungen} & {\leavevmode\small\color{Muted}#1 von 3} & {\leavevmode\small\color{Muted}Lösungsblatt \textperiodcentered{} Aufgabe 7}\end{tabularx}}
\tcbset{sheetbox/.style={colframe=Line,arc=1.2mm,boxrule=.6pt,left=2mm,right=2mm,top=1mm,bottom=1mm,before skip=1.5mm,after skip=1.5mm}}
\setlength{\abovedisplayskip}{4pt}
\setlength{\belowdisplayskip}{4pt}
\begin{document}
\sheethead{1 \textperiodcentered{} Hammerwurf}
\begin{tcolorbox}[sheetbox,colback=Sky]
\textbf{Gegeben:} $l_{\mathrm{Seil}}=1{,}22\unit{m}$, $l_{\mathrm{Arm}}=86{,}0\unit{cm}$,
$v_{\mathrm B}=25{,}0\unitfrac{m}{s}$, $m=7{,}16\unit{kg}$.\\[1mm]
Gesucht: Frequenz $f$, Haltekraft $F_{\mathrm Z}$ und geeignete Abwurfposition.
\end{tcolorbox}
\sectiontitle{1a \textperiodcentered{} Frequenz kurz vor dem Abwurf}
\step{1}{Armlänge umrechnen und Radius bestimmen}
Der Radius reicht von der Drehachse des Werfers bis zur Kugel. Deshalb werden Armlänge und Seillänge addiert:
\[
l_{\mathrm{Arm}}=86{,}0\unit{cm}=0{,}860\unit{m},\qquad
r=1{,}22\unit{m}+0{,}860\unit{m}=2{,}08\unit{m}.
\]
\step{2}{Zusammenhang aufstellen und nach der Frequenz auflösen}
In einer Umdrehung legt die Kugel den Umfang $2\pi r$ zurück. Mit $f$ Umdrehungen pro Sekunde gilt:
\[
v_{\mathrm B}=2\pi\cdot f\cdot r
\quad\Longrightarrow\quad f=\frac{v_{\mathrm B}}{2\pi\cdot r}.
\]
\step{3}{Einsetzen und auf drei gültige Ziffern runden}
\[
f=\frac{25{,}0\unitfrac{m}{s}}{2\pi\cdot2{,}08\unit{m}}
=1{,}9129\ldots\unit{Hz}\quad\Longrightarrow\quad\boxed{f\approx1{,}91\unit{Hz}}.
\]
Die Kugel führt kurz vor dem Abwurf etwa $1{,}91$ Umdrehungen pro Sekunde aus. Die drei vorherigen Drehungen werden für diese momentane Frequenz nicht benötigt.
\sectiontitle{1b \textperiodcentered{} Haltekraft kurz vor dem Abwurf}
\step{1}{Physikalischen Ansatz wählen}
Im vereinfachten Modell der Aufgabe ist die gesuchte Haltekraft die zum Kreismittelpunkt gerichtete Zentripetalkraft:
\[
F_{\mathrm Z}=\frac{m\cdot v_{\mathrm B}^{2}}{r}.
\]
\step{2}{Werte einsetzen; die Geschwindigkeit quadrieren}
\[
F_{\mathrm Z}=\frac{7{,}16\unit{kg}\cdot\left(25{,}0\unitfrac{m}{s}\right)^2}{2{,}08\unit{m}}
=2151{,}44\ldots\unit{N}.
\]
\step{3}{Mit drei gültigen Ziffern angeben}
\begin{tcolorbox}[sheetbox,colback=Mint]
\centering $\boxed{F_{\mathrm Z}\approx2{,}15\cdot10^3\unit{N}=2{,}15\unit{kN}}$\quad Die Kraft zeigt zur Drehachse.
\end{tcolorbox}
\sectiontitle{1c \textperiodcentered{} Geeignete Stelle zum Loslassen}
\begin{minipage}[c]{.64\linewidth}
\textbf{1. Bewegungsrichtung bestimmen:} Die momentane Geschwindigkeit zeigt tangential zur Kreisbahn.\\[1.5mm]
\textbf{2. Wirkung des Loslassens erklären:} Die Kraft zur Drehachse entfällt. Von oben gesehen bewegt sich der Hammer tangential weiter.\\[1.5mm]
\textbf{3. Abwurfposition wählen:} Der Hammer muss dort losgelassen werden, wo die Tangente in Bewegungsrichtung durch die Netzöffnung ins freie Feld zeigt. In der Skizze ist das für den dargestellten Drehsinn der rote Punkt.
\end{minipage}\hfill
\begin{minipage}[c]{.33\linewidth}\centering
\begin{tikzpicture}[x=1mm,y=1mm,>=Latex]
\coordinate (W) at (16,16);
\draw[Muted,line width=1.6pt] ($(W)+(50:14)$) arc[start angle=50,end angle=310,radius=14];
\node[font=\scriptsize,Muted,anchor=south west] at ($(W)+(55:14.5)$) {Schutznetz};
\draw[Blue,line width=.7pt] (W) circle (9);
\draw[-{Latex},Blue,line width=.7pt] ($(W)+(155:9)$) arc[start angle=155,end angle=105,radius=9];
\fill[Navy] (W) circle (1.2);
\node[font=\scriptsize,below] at ($(W)+(0,-1.4)$) {Werfer};
\draw[-{Latex},green!50!black,line width=1.1pt] ($(W)+(0,9)$) -- ($(W)+(31,9)$);
\fill[red!80!black] ($(W)+(0,9)$) circle (1.3);
\end{tikzpicture}
\par\vspace{1mm}{\scriptsize\color{Muted}Blick von oben, nicht ma\ss{}stäblich. Die Gewichtskraft wirkt weiterhin nach unten.}
\end{minipage}
\pagefooter{1}
\newpage
\sheethead{2 \textperiodcentered{} Rennbahn}
\sectiontitle{2 \textperiodcentered{} Optimale Neigung einer Rennbahn}
\begin{tcolorbox}[sheetbox,colback=Sky]
\textbf{Gegeben:} $r=50\unit{m}$, $v_{\mathrm B}=36\unitfrac{km}{h}$, $g=9{,}81\unitfrac{m}{s^2}$.\\[1mm]
\textbf{Gesucht:} Neigung $\alpha$ gegen die Waagerechte, bei der keine seitliche Reibung nötig ist.
\end{tcolorbox}
\step{1}{Geschwindigkeit in SI-Einheiten umrechnen}
\[
v_{\mathrm B}=36\unitfrac{km}{h}=\frac{36\cdot1000\unit{m}}{3600\unit{s}}=10\unitfrac{m}{s}.
\]
\step{2}{Kräfte und Winkel im Kräftedreieck zuordnen}
Es wirken die Gewichtskraft $F_{\mathrm G}$ senkrecht nach unten und die Normalkraft $F_{\mathrm N}$ senkrecht zur Bahn. Ihre Summe ist waagerecht zur Kurvenmitte gerichtet:
\[
\vec F_{\mathrm N}+\vec F_{\mathrm G}=\vec F_{\mathrm Z}.
\]
\begin{minipage}[c]{.40\linewidth}\centering
\begin{tikzpicture}[x=.95mm,y=.95mm,>=Latex]
\coordinate (S) at (0,0);\coordinate (P) at (16,34.3);\coordinate (Q) at (16,0);
\draw[-{Latex},Blue,line width=1.2pt] (S)--(P) node[midway,left] {$F_{\mathrm N}$};
\draw[-{Latex},Gold!85!black,line width=1.2pt] (P)--(Q) node[midway,right] {$F_{\mathrm G}$};
\draw[-{Latex},red!80!black,line width=1.2pt] (S)--(Q) node[midway,below] {$F_{\mathrm Z}$};
\draw ($(Q)+(-3,0)$)--++(0,3)--++(3,0);
\draw ($(P)+(270:9)$) arc[start angle=270,end angle=245,radius=9];
\node at ($(P)+(257:12)$) {$\alpha$};
\end{tikzpicture}
\par\vspace{1mm}{\scriptsize\color{Muted}Kräftedreieck, nicht ma\ss{}stäblich.}
\end{minipage}\hfill
\begin{minipage}[c]{.57\linewidth}
$F_{\mathrm N}$ ist die \textbf{Hypotenuse}, $F_{\mathrm G}$ die \textbf{Ankathete} und $F_{\mathrm Z}$ die \textbf{Gegenkathete} von $\alpha$.\\[2mm]
Da $F_{\mathrm N}$ senkrecht zur Bahn steht, ist der Winkel zur Senkrechten genauso gro\ss{} wie die Neigung der Bahn gegen die Waagerechte.
\end{minipage}
\step{3}{Senkrechte und waagerechte Teilgleichung aufstellen}
Die Höhe bleibt gleich, deshalb heben sich die senkrechten Kraftanteile auf. Waagerecht ist die Zentripetalkraft nötig:
\begin{align*}
F_{\mathrm N}\cdot\cos\alpha&=F_{\mathrm G}=m\cdot g,\\
F_{\mathrm N}\cdot\sin\alpha&=F_{\mathrm Z}=\frac{m\cdot v_{\mathrm B}^{2}}{r}.
\end{align*}
\step{4}{Waagerechte durch senkrechte Gleichung teilen und kürzen}
\[
\frac{F_{\mathrm N}\cdot\sin\alpha}{F_{\mathrm N}\cdot\cos\alpha}
=\frac{m\cdot v_{\mathrm B}^{2}}{r\cdot m\cdot g}
\quad\Longrightarrow\quad
\frac{\sin\alpha}{\cos\alpha}=\frac{v_{\mathrm B}^{2}}{r\cdot g}.
\]
Links kürzt sich $F_{\mathrm N}$, rechts $m$. Mit $\dfrac{\sin\alpha}{\cos\alpha}=\tan\alpha$ folgt:
\[
\tan\alpha=\frac{v_{\mathrm B}^{2}}{r\cdot g}.
\]
\step{5}{Einsetzen und mit der Umkehrfunktion den Winkel berechnen}
\[
\tan\alpha=\frac{\left(10\unitfrac{m}{s}\right)^2}{50\unit{m}\cdot9{,}81\unitfrac{m}{s^2}}
=0{,}203873\ldots
\]
\[
\alpha=\arctan\left(\frac{\left(10\unitfrac{m}{s}\right)^2}{50\unit{m}\cdot9{,}81\unitfrac{m}{s^2}}\right)
=11{,}523\ldots{}^{\circ}.
\]
\begin{tcolorbox}[sheetbox,colback=Mint]
\textbf{Ergebnis mit zwei gültigen Ziffern:} $\boxed{\alpha\approx12^{\circ}}$. Die Bahn muss zur Kurvenmitte hin abfallen. Die optimale Neigung hängt nicht von der Masse des Fahrzeugs ab.
\end{tcolorbox}
\begin{tcolorbox}[sheetbox,colback=Sky]
\textbf{Beim Rechnen beachten:} Taschenrechner auf Grad (DEG) stellen. Erst das Endergebnis runden. $\tan\alpha\approx0{,}204$ ist ein Verhältnis, noch kein Winkel.
\end{tcolorbox}
\pagefooter{2}
\newpage
\sheethead{3 \textperiodcentered{} Kettenkarussell und Abschlussquiz}
\small\setlength{\abovedisplayskip}{2pt}\setlength{\belowdisplayskip}{2pt}
\renewcommand{\step}[2]{\par\vspace{1mm}\textbf{\textcolor{Blue}{#1}\quad #2}\par\vspace{.3mm}}
\sectiontitle{3a \textperiodcentered{} Kräfte am Mann einzeichnen}
\begin{minipage}[c]{.31\linewidth}\centering
\begin{tikzpicture}[x=.58mm,y=.58mm,>=Latex]
\coordinate (A) at (6,45);\coordinate (P) at (46,18);
\coordinate (FG) at (46,-2);\coordinate (FZ) at (16.3,18);\coordinate (FS) at (16.3,38);
\draw[Muted,dash dot,line width=.6pt] (6,-4)--(6,54);
\node[font=\scriptsize,anchor=south] at (6,54) {Drehachse};
\draw[Muted,line width=.7pt] (A)--(P) node[pos=.17,above right,font=\scriptsize] {$l$};
\draw[Muted,{Latex[length=1.6mm]}-{Latex[length=1.6mm]}] (6,50)--(46,50) node[midway,above,font=\scriptsize] {$r$};
\draw[Muted,densely dotted] (46,50)--(46,21);
\draw ($(A)+(-90:8)$) arc[start angle=-90,end angle=-34,radius=8];
\node[font=\scriptsize] at (9.5,33.5) {$\alpha$};
\draw[Blue!45,densely dashed] (FS)--(FZ);\draw[Blue!45,densely dashed] (FG)--(FZ);
\fill (P) circle (.9);
\draw[-{Latex},Gold!85!black,line width=1.1pt] (P)--(FG) node[right,font=\scriptsize] {$F_{\mathrm G}$};
\draw[-{Latex},Blue,line width=1.1pt] (P)--(FS) node[pos=.55,below left,font=\scriptsize] {$F_{\mathrm S}$};
\draw[-{Latex},red!80!black,line width=1.1pt] (P)--(FZ) node[midway,below,font=\scriptsize] {$F_{\mathrm Z}$};
\end{tikzpicture}
\par{\scriptsize\color{Muted}Seitenansicht, nicht ma\ss{}stäblich.}
\end{minipage}\hfill
\begin{minipage}[c]{.66\linewidth}
\textbf{1.} Am Körperschwerpunkt $F_{\mathrm G}$ senkrecht nach unten einzeichnen.\\[1mm]
\textbf{2.} $F_{\mathrm S}$ entlang der Kette zum Aufhängepunkt einzeichnen. Die Seilkraft ist grö\ss{}er als die Gewichtskraft, denn nur ihr senkrechter Anteil gleicht $F_{\mathrm G}$ aus.\\[1mm]
\textbf{3.} Das Kräfteparallelogramm ergänzen. Seine Diagonale $F_{\mathrm Z}$ zeigt waagerecht zur Drehachse:
\[
\vec F_{\mathrm S}+\vec F_{\mathrm G}=\vec F_{\mathrm Z}.
\]
$F_{\mathrm Z}$ bezeichnet die \textbf{Resultierende}, keine dritte äu\ss{}ere Kraft.
\end{minipage}
\sectiontitle{3b \textperiodcentered{} Bahngeschwindigkeit berechnen}
\textbf{Gegeben:} $l=15\unit{m}$, $\alpha=56^{\circ}$, $m=75\unit{kg}$, $g=9{,}81\unitfrac{m}{s^2}$.
\step{1}{Radius aus der Geometrie bestimmen}
Der Aufhängepunkt liegt im Aufgabenmodell an der Drehachse. Im rechtwinkligen Dreieck ist $l$ die Hypotenuse und $r$ die Gegenkathete zu $\alpha$:
\[
\sin\alpha=\frac{r}{l}\quad\Longrightarrow\quad
r=l\cdot\sin\alpha=15\unit{m}\cdot\sin56^{\circ}=12{,}43556\ldots\unit{m}.
\]
\step{2}{Kräftegleichungen aufstellen und durcheinander teilen}
\[
F_{\mathrm S}\cdot\cos\alpha=m\cdot g,\qquad
F_{\mathrm S}\cdot\sin\alpha=\frac{m\cdot v_{\mathrm B}^{2}}{r}.
\]
\[
\frac{F_{\mathrm S}\cdot\sin\alpha}{F_{\mathrm S}\cdot\cos\alpha}
=\frac{m\cdot v_{\mathrm B}^{2}}{r\cdot m\cdot g}
\quad\Longrightarrow\quad\tan\alpha=\frac{v_{\mathrm B}^{2}}{r\cdot g}.
\]
Die Seilkraft und die Masse kürzen sich. Die Angabe $m=75\unit{kg}$ wird für die Geschwindigkeit nicht benötigt.
\step{3}{Nach der Bahngeschwindigkeit auflösen}
\[
v_{\mathrm B}^{2}=r\cdot g\cdot\tan\alpha
\quad\Longrightarrow\quad v_{\mathrm B}=\sqrt{r\cdot g\cdot\tan\alpha}.
\]
\step{4}{Mit ungerundetem Radius einsetzen und am Ende runden}
\[
v_{\mathrm B}=\sqrt{15\unit{m}\cdot\sin56^{\circ}\cdot9{,}81\unitfrac{m}{s^2}\cdot\tan56^{\circ}}
=13{,}4485\ldots\unitfrac{m}{s}.
\]
\begin{tcolorbox}[sheetbox,colback=Mint]
\textbf{Zwei gültige Ziffern:} $\boxed{v_{\mathrm B}\approx13\unitfrac{m}{s}}$. Aus dem ungerundeten Wert folgt alternativ $v_{\mathrm B}=48{,}414\ldots\unitfrac{km}{h}\approx48\unitfrac{km}{h}$.
\end{tcolorbox}
\sectiontitle{4 \textperiodcentered{} Abschlussquiz: richtige Aussagen und Begründungen}
\small
Die Nummern nennen die richtigen Auswahlkästchen in der Reihenfolge der Aussagen im Modul.
\par\vspace{1mm}
\begin{tabularx}{\linewidth}{@{}>{\bfseries}l >{\raggedright\arraybackslash}p{19mm} >{\raggedright\arraybackslash}X@{}}
\toprule
Frage & Kästchen & Begründung\\\midrule
4a & 1 und 3 & Die Zentripetalkraft ist die zum Kreismittelpunkt gerichtete Resultierende realer Kräfte.\\[1mm]
4b & 1 und 3 & Doppelte Masse oder halber Radius verdoppeln $F_{\mathrm Z}$. Doppelte Geschwindigkeit vervierfacht $F_{\mathrm Z}$; doppelter Radius halbiert sie.\\[1mm]
4c & 1 und 3 & $F_{\mathrm N}\cos\alpha=F_{\mathrm G}$ und $F_{\mathrm N}\sin\alpha=F_{\mathrm Z}$; $F_{\mathrm N}$ ist die Hypotenuse.\\[1mm]
4d & 1 und 3 & Die Masse kürzt sich. Bei festem Radius verlangt eine höhere Geschwindigkeit eine grö\ss{}ere Neigung; ein grö\ss{}erer Radius verlangt bei gleicher Geschwindigkeit eine kleinere.\\[1mm]
4e & 2 und 4 & Der Mann spürt seine Trägheit. Rei\ss{}t die Kette, bewegt er sich von oben gesehen tangential weiter. Für den Beobachter am Boden wirkt keine Fliehkraft nach au\ss{}en.\\[1mm]
4f & 1, 3 und 5 & Im Modell wirken $F_{\mathrm G}$ und $F_{\mathrm S}$. Ihre Summe zeigt zur Drehachse und ist nicht null. Wegen $F_{\mathrm S}\cos\alpha=F_{\mathrm G}$ gilt $F_{\mathrm S}>F_{\mathrm G}$.\\
\bottomrule
\end{tabularx}
\pagefooter{3}
\end{document}
